\documentclass[12pt]{article}
\usepackage{amsmath,amssymb,amsfonts}
\usepackage[margin=1in]{geometry}

\begin{document}

\section*{Chapter 3, Problem 13}

\textbf{Problem.} Prove that $\operatorname{tr}(AB) = \operatorname{tr}(BA)$ for arbitrary $n \times n$ matrices $A$ and $B$. Use this result to prove that if $C$ and $D$ are similar, then $\operatorname{tr}(C) = \operatorname{tr}(D)$.

\bigskip
\textbf{Solution.}

\textbf{Part 1:} $\operatorname{tr}(AB) = \operatorname{tr}(BA)$.

By definition of the trace:
\[
\operatorname{tr}(AB) = \sum_{i=1}^{n} (AB)_{ii} = \sum_{i=1}^{n} \sum_{j=1}^{n} A_{ij} B_{ji}.
\]
Similarly:
\[
\operatorname{tr}(BA) = \sum_{j=1}^{n} (BA)_{jj} = \sum_{j=1}^{n} \sum_{i=1}^{n} B_{ji} A_{ij}.
\]
Since these are finite double sums of the same terms $A_{ij}B_{ji}$ (just with the order of summation interchanged), they are equal:
\[
\operatorname{tr}(AB) = \sum_{i,j} A_{ij}B_{ji} = \operatorname{tr}(BA). \quad \checkmark
\]

\medskip
\textbf{Part 2:} If $C$ and $D$ are similar, i.e., $D = P^{-1}CP$ for some invertible matrix $P$, then:
\[
\operatorname{tr}(D) = \operatorname{tr}(P^{-1}CP).
\]
Using the cyclic property of the trace (which follows from Part 1), let $A = P^{-1}C$ and $B = P$:
\[
\operatorname{tr}(P^{-1}CP) = \operatorname{tr}(P \cdot P^{-1}C) = \operatorname{tr}(C).
\]

Therefore $\operatorname{tr}(D) = \operatorname{tr}(C)$. The trace is invariant under similarity transformations. \hfill $\blacksquare$

\end{document}
